\subsectionpaged{Calculate From}
\label{config:calculateFrom}
\index{calculate from}
\index{Config!calculate from}

\begin{lstlisting}[language=Go,style=kaolstplain,escapechar=|]
calculate from "metric"
    [unit "unit"]
    ( |\textit{aggregates}\dots| )
    [reset every "crontab"]
\end{lstlisting}

\Statement{calculate from} calculates one or more aggregate metrics from a single source metric.

It is a convenience form of \Statement{calculate()} removing a lot of boilerplate configuration into a single statement.

\ParameterDescription{
\item[unit]defines the unit we want the results to be in.
If not defined then the unit of the output of the aggregate function is used.

\item[aggregates]define the aggregations to perform. The available functions are listed in \vref{config:calculatefor:aggregate:table}.

\item[reset every]defines when the metric being calculated should be reset to some default value.
}

\begin{table}[ht]
    \caption{Available aggregate functions for calculate for}
    \label{config:calculatefor:aggregate:table}
    \begin{tabular}{l l}
        \toprule
        Aggregator & Function \\
        \midrule
        avg & Average \\
        min & Minimum \\
        max & Maximum \\
        sum & Summation \\
        \bottomrule
    \end{tabular}
\end{table}

\subsubsectionfit{Example}
The following example shows how to calculate the minimum and maximum values of a specific metric.

\begin{lstlisting}[language=Go,caption={Example calculation for},escapechar=|,label=lst:calculatefrom:example]
calculate from "ws90.humidity"  |\label{line:calculatefrom:example:src}|
    unit "relativeHumidity"     |\label{line:calculatefrom:example:unit}|
    ( avg min max )             |\label{line:calculatefrom:example:aggregators}|
    reset every "day"           |\label{line:calculatefrom:example:reset}|
\end{lstlisting}

Line \ref{line:calculatefrom:example:src} declares the metric we will trigger the calculations from.

Line \ref{line:calculatefrom:example:unit} defines the unit we want the results to be recorded under.

Line \ref{line:calculatefrom:example:aggregators} declares which aggregators to use.

Line \ref{line:calculatefrom:example:reset} indicates we want the metrics to be reset once every day at midnight.

The result of this is that three new metrics are generated using \Metric{ws90.humidity} as the source.
As we are using \Statement{avg}, \Statement{min} and \Statement{max} here then this will create the metrics
\Metric{ws90.humidity.avg}, \Metric{ws90.humidity.min} and \Metric{ws90.humidity.max} respectively.

\subsubsectionfit{Expanding the example}
The configuration in \vref{lst:calculatefrom:example} is equivalent to the following standard configuration:

\begin{lstlisting}[language=Go,caption={calculation for example expanded},escapechar=|,label=lst:calculatefrom:expanded]
calculate( "ws90.humidity.avg"
    reset every "day"
    load "today" with "avg(ws90.humidity)"
    usefirst "ws90.humidity"
    as avg(current,"ws90.humidity") unit "relativeHumidity"
)

calculate( "ws90.humidity.min"
    reset every "day"
    load "today" with "min(ws90.humidity)"
    usefirst "ws90.humidity"
    as min(current,"ws90.humidity") unit "relativeHumidity"
)

calculate( "ws90.humidity.max"
    reset every "day"
    load "today" with "max(ws90.humidity)"
    usefirst "ws90.humidity"
    as max(current,"ws90.humidity") unit "relativeHumidity"
)

\end{lstlisting}
